%\documentclass[handout]{beamer}
\documentclass{beamer}
% \usetheme{Rochester}
\usecolortheme{orchid}
\def\ol{\overline}
\usepackage{amsthm}
\newtheorem{prop}{Proposal}
\newtheorem{ansatz}{Ansatz}
\newtheorem{thm}{Theorem}
\title{Antimatter and the Metaphysics of Quantity} 
\author{Hans Halvorson \\ (with David Baker) \\
  {\small Princeton University}} 
\date{November 30, 2007}
\begin{document}

\begin{frame}
\titlepage
\end{frame}

%% Is antimatter discovery trivial ... like discovering that there is
%% another car going exactly the same speed, but to the south?


% \section*{Introduction}
% \begin{frame}
% \begin{itemize}
% \item Basic concepts of contemporary metaphysics: object, property,
%   relation, disposition, time, etc..  
% % Can also define quantities, etc., in terms of basic concepts.

% \item The notion of \textbf{opposite} is not ever invoked ever in
%   contemporary metaphysics, and is certainly not thought to be
%   fundamental.

% \item Today's goals: 
% \begin{itemize}
% \item Use contemporary physics to motivate and explicate a concept of
%   ``opposite''.
% \item Show that the opposite concept is fundamental.
% \end{itemize}
% \end{itemize}
% \end{frame}

% \begin{frame}
% \frametitle{Metaphysical Lessons}

% %% MOTIVATION: WHY SHOULD PHILOSOPHERS CARE %%

% %% asymmetry of matter vs. antimatter
% %% what does this tell us?
% %% is it connected to assymmetry of time?

% %% two concepts of "negation" - semantic and ontological??

% %% lessons for metaphysics: contemporary analytic metaphysics seems to
% %% have no use for a concept of the *opposite* of a thing

% %% 20th century physics, reistatement of the concept of "vacuum" 

% %% What is vacuum? It is nothing.

% %%  
% \pause
% \begin{enumerate}
% \item Each essential quantity $\mathcal{C}$ has a unique privileged
%   neutral value $1_{\mathcal{C}}$.
% \item There is a relation of ``opposition'' between fundamental
%   properties.
%   %% not definable in terms of other physical relations; and applies
%   %% to each fundamental physical object. 
% %% note that this relation is symmetric and it *can* hold between an
%   %% object and itself
% \item Metaphysical accounts of \textbf{quantity} should be revised in
%   light of recent empirical discoveries.
% \end{enumerate}


% \end{frame}

\section*{Outline}
\begin{frame}
\frametitle{Outline}
\tableofcontents
\end{frame}


\section{Philosophical accounts of quantity}
\begin{frame}
  \frametitle{Quantities and determinables}
\pause
\begin{tabular}{cc|cc}
  Determinable && Determinate \\ 
  \hline\hline & & & \\ color && red, green, blue, \dots \\ &&& \\
  \hline &&& \\ mass && non-negative real number \\ &&& \\
  \hline &&& \\ position &&  point in space  \\ &&& \\
  \hline &&& \\ isospin  &&  $0\,,\,\frac{1}{2}\,,\,1\,,\,\frac{3}{2}\,,\,\dots$ 
\end{tabular}


\end{frame}
\begin{frame}
  \begin{prop} Quantities and determinables are the same
    thing.  \end{prop} \pause
  \begin{prop} A determinable/quantity is a collection of mutually
    exclusive and exhaustive properties.  \end{prop} \pause
%% note to self: so, this  would mean that having the value zero is
%% having a definite value ... not lacking any determinate value for
%% that determinable
  \begin{prop}[Mundy]
    A quantity is a semi-group; i.e.\ a set $S$ with a binary function
    $\quad\star :S\times S\to S$.
\end{prop} \pause
\begin{prop}[Denby] 
  A quantity is a pair $\langle S,f\rangle$ where $S$ is a
  \emph{metric space}, and $f$ is a function from concrete particulars
  to $S$. \end{prop}
\end{frame}


\begin{frame}\frametitle{Causes for worry}
\pause

\begin{itemize}
\item An account of the metaphysics of quantity should especially
  accommodate the use of the concept in fundamental physics. \pause
\item Mundy's account is based on examples from phenomenological
  physics (measurement theory). \pause
\item Denby's account is based on a pre-scientific notion of
  ``similarity'' between properties --- e.g.\ ``redness'' is more
  similar to ``purpleness'' than it is to ``blueness''.
\end{itemize}
\end{frame}


\begin{frame}
\frametitle{Input from current physics?} \pause
\begin{itemize}
\item Fine and Teller (circa 1978): Continuous quantities have only
  fuzzy values.  \pause Halvorson (2001): No. \pause Arntzenius
  (2003): Yes. \pause
\end{itemize}
\begin{itemize}
\item Determinates are structured \dots but \emph{not} necessarily
  with the structure of a semigroup, or a metric.  \pause
  %% Here I'm thinking of isospin, but .... hmm .... wouldn't we say
  %% that isospin 1 is closer to 1/2 than is 2?

  %% Is there a natural topology on the category of representations of
  %% a compact group?  Is it always a metric topology?
  %% Look at the Folland book -- apparently it is always the discrete
  %% topology
\end{itemize}
\end{frame}

\section{A brief history of antimatter}
\begin{frame}
\frametitle{A brief history of antimatter}
\begin{itemize}
\item 1927: Dirac's equation for relativistic electrons permits two
  solutions, one with positive energy and one with negative energy. \pause
\item 1932: Carl Anderson confirms the existence of the antielectron
  (positron); 1955: antiproton. \pause
  %% antiparticle of the electron
\item ca 1970--80: The standard model entails that every fundamental
  particle has an antiparticle. \pause
\item 1996: Synthesis of atoms of antihydrogen (antiproton+positron)
  at CERN. \pause
\item Ongoing: Apparent asymmetry of matter and antimatter in our
  universe is one of the greatest unsolved problems in contemporary
  physics. \pause
\item 1969--1990: Doplicher, Haag, and Roberts explicate the structure
  of fundamental physical quantities
\end{itemize}

\end{frame}



\section{What antimatter could not be}
\begin{frame}
\frametitle{Is antimatter a \emph{philosophical} topic?}
\pause

Two extreme, but wrong, positions:
\begin{itemize}
\item Contemporary physics shows that besides ``being'' (or existence)
  there is ``anti-being'' (or anti-existence).  The intermediate
  between the two is nothingness (the void). \pause
\item ``Antimatter'' is a misnomer: everything is matter. \pause
\end{itemize}


\bigskip Other reason to care about antimatter: Asymmetry of
matter/antimatter connected to asymmetry of future/past? 



\end{frame}

\begin{frame}

\begin{prop} The antiparticle $\ol p$ of particle $p$ is something
    that is disposed to \emph{annihilate}~$p$.  \end{prop}

\bigskip \pause
\begin{itemize}
\item Not a good analysis (cf.\ ``opium causes sleep because it has a
  dormative virtue'') \pause
\item Antimatter is definable even in theories that have no
  interactions. \pause
\item Not a necessary condition: Antiparticles do not always
  annihilate their antiparticles, but sometimes produce other
  particles \pause
\item Not a necessary condition: Some particles are their own
  antiparticles (e.g.\ photon, neutral pion)
\end{itemize}
\end{frame}


\begin{frame}
  \begin{prop}[Feynman, Arntzenius and Greaves] Antiparticles are
    particles moving backwards in time. \end{prop} \pause

  \begin{itemize}
  \item Feynman's heuristic reasoning does not amount to an
    analysis. \pause
  \item This claim can be explicated as a mathematical statement that
    is provably true \dots but which undercuts the provocative
    interpretation of particles moving in opposite directions in
    time. \pause
  \item The claim is not general enough: the antimatter concept is
    more general than the particle concept.

  \end{itemize}

\end{frame}


\begin{frame}
  \begin{prop} Antiparticles have conjugate wavefunctions. \end{prop}
  \bigskip \pause

%% this would undercut fundamentality, because it would make the
%% antimatter concept be contingent on the idiosyncratic fact that our
%% best theory (viz.\ QM) uses complex amplitudes

  \textbf{Definition:} A \emph{wavefunction} $\psi$ is an assignment
  of complex numbers to the determinates of a determinable.

  \bigskip A (nonzero) complex number is a pair consisting of a
  length, and an angle.
$$ \psi (e) \quad = \quad (\text{probability},\text{phase}) $$

\bigskip Complex Conjugation
$$(\text{probability},\text{phase}) \quad
\stackrel{\text{conjugate}}{\Longrightarrow} \quad
(\text{probability},-\text{phase})$$

\end{frame}

\begin{frame}
\frametitle{Conjugate wavefunction}

\begin{itemize}
\item Definition:
$$\overline{\psi}(e) \quad =\quad (\psi (e))^* $$
The wavefunctions $\psi$and $\ol\psi$ assign the same probabilities to
all determinate values. \pause
\item Wrong: the conjugate wavefunction is basis-dependent; cannot
  ground a fundamental physical relation.
%% similar to problem with setting zero arbitrarily
\end{itemize}

\end{frame}


\begin{frame}
  \begin{prop}[Oxford Dictionary of Physics] The antiparticle of $p$
    is the particle $\ol p$ with \emph{opposite} values for each
    quantity. \end{prop}


\bigskip \pause
\begin{itemize}
\item Strictly false: some quantities' values are not reversed (e.g.\
  position, velocity) \pause
\item Even false for ``essential'' quantities.  (e.g.\ mass) \pause
  %% e.g. mass of electron and positron are the same e.g. isospin of
  %% ...
\item The concept of ``opposite value'' is left unanalyzed \newline On
  going philosophical accounts of quantity (e.g.\ Mundy) the specific
  assignment of real numbers to properties is conventional
\end{itemize}
%% additive quantum numbers??
\end{frame}


\begin{frame}
  \begin{prop}[Modified dictionary definition; Penrose] The
    antiparticle of $p$ is the particle $\ol p$ with \emph{opposite}
    values for each \emph{additive} essential quantity. \end{prop}

  %% quote from Oxford dictionary
\bigskip \pause

\begin{quote}``\ldots a quantum number \dots is `additive' if, in
  order to derive its value for a composite entity, we simply add up
  the individual values for the constituent particles'' \end{quote}

\end{frame}






\begin{frame}
  \frametitle{The opposite of a property}

\bigskip {\small Newton's third law: For every force acting on an
  object, the object will exert an equal, yet \emph{opposite}, force
  on its cause.}  \bigskip \bigskip

\begin{center}
\begin{tabular}{ll|ll}
  Opposites  & & & No Opposites \\
  \hline 
   & & & Temperature \\
  Newtonian Forces & & &   \\
  Electric Charge & & &  Position in post-Aristotelian space \\
  Relative Position in STR &&& Relative Position in GTR  \\
  &&& Ounces of gold     \\
  Essential quantities in &&& \\ 
  quantum field theory \end{tabular}
\end{center}
\end{frame}




\section{Antimatter and the structure of quantities}

\begin{frame}
  \frametitle{Abelian quantum numbers (e.g.\ electric charge)}
% isn't there a common theme here in that they are both "forces"?
% isn't GTR based on the idea that forces themselves are not
% absolute?
% cf. Norton's article about Newtonian cosmology
\begin{itemize}
\item There are an infinite number of charges $\{ \dots
  ,-2,-1,0,1,2,\dots \}$.
\item Charges can be \textbf{composed}, and the composition operation
  is modelled by addition of real numbers.  
\item There is a unique \textbf{neutral} charge, to which we assign
  value $0$.
% \item The gauge transformations of electric charge are also a group.
% \item These two groups are ``duals'' of each other.
\end{itemize}
\end{frame}



\begin{frame}
  \frametitle{Nonabelian quantum numbers (e.g.\ isospin)}
Isospin takes values in $\{ 0,\frac{1}{2},1,\frac{3}{2},\dots \}$

Composition rules:
\begin{eqnarray*} 
i\otimes 0 &=& 0\otimes i \:=\: i \\ \\
\textstyle \frac{1}{2}\otimes \frac{1}{2}& =& 0\oplus 1 \\ \\
\textstyle 1\otimes \frac{1}{2}&=& \textstyle \frac{1}{2}\oplus 1 \oplus
\frac{3}{2} \\ \\
i\otimes j &= & \bigoplus _{k=|i-j|}^{i+j}k \end{eqnarray*}
\end{frame}


\begin{frame}
\frametitle{Nonabelian quantum numbers}
What is this second operation ``$\oplus$'' that appears on the right
hand side of equations involving ``$\otimes$''? \pause
\begin{itemize} 
\item Role in empirical predictions: \newline \textit{If $i\otimes
    j=k_1\oplus \cdots \oplus k_n$, then a collision of particles with
    quantum numbers $i$ and $j$ results in a mixture of particles with
    quantum numbers $k_1,\dots ,k_n$.} \pause
\item Mathematical structure: \newline $x=k_1\oplus \cdots \oplus k_n$
  means that $k_1,\dots ,k_n$ are ``parts'' of $x$, i.e.\ there are
  orthogonal embeddings of the objects $k$ in the object $x$.  (Each
  of the objects $0,\frac{1}{2},1,\frac{3}{2},\dots$ is itself
  irreducible.) \pause
\end{itemize}

\end{frame}

\begin{frame}
\frametitle{Conjugate objects}
\textbf{Definition:} The conjugate of $i$ is the unique object $j$
such that $0$ occurs in the decomposition of $i\otimes j$. \pause

\bigskip
\begin{itemize}
\item Electric charge: the conjugate of $i$ is $-i$, and $i\neq
  -i$.\pause
\item Isospin: the conjugate of $i$ is $i$.
\end{itemize}
\end{frame}


\begin{frame}\frametitle{The magic symmetry principle}

  Gell-Mann et al.: The quantum numbers are objects of the dual of the
  symmetry group of the theory. \pause

  \bigskip

  \textbf{Examples:} 
\begin{itemize}
\item The symmetry group for isospin is $SU(2)$.  In the dual of
  $SU(2)$, each object is self-conjugate; hence, each isospin quantum
  number is its own inverse. \pause 
\item The symmetry group for charge is $U(1)$, and the dual of
  $U(1)$ is the group $\mathbb{Z}$ of integers. \pause \end{itemize}

\bigskip

\textbf{Problem:} Why does this work?  What does the dual of the
  symmetry group have to do with physical quantities?

\end{frame}


\begin{frame}
  \frametitle{The algebraic approach to quantum field theory}

\begin{itemize}
\item von Neumann: Schr{\"o}dinger's wave mechanics and Heisenberg's
  matrix mechanics are provably equivalent. \pause
\item van Hove, Friedrichs, Haag: Because quantum field theory (QFT)
  has infinitely many degrees of freedom, von Neumann's uniqueness
  theorem fails --- there are unitarily inequivalent
  representations. \pause
\item The existence of UIRs has been an interpretive stumbling block
  for philosophers. \pause
\end{itemize}

\begin{ansatz}[Doplicher, Haag, Roberts]
  Each UIR corresponds to an essential determinate (value of a
  quantity) of elementary particles.
\end{ansatz} 
\end{frame}

\begin{frame}
  \frametitle{DHR program: derive structure of quantities}

  \textbf{Question:} Does the ``set'' of UIRs have an interesting
  additional structure --- derivable from the physics? \pause

\bigskip 
\textbf{Answer:} Yes.  For any two representations $\pi$ and $\pi '$:
\begin{itemize}
\item There is a product representation $\pi\otimes \pi '$;
\item There is a sum representation $\pi \oplus \pi '$;
\item There is a unique identity representation $\pi _0$ such that
  $\pi \otimes \pi _0 = \pi =\pi _0\otimes \pi$;
\item There is a unique conjugate representation $\ol\pi$ such that
  $\pi\otimes \ol\pi$ and $\ol\pi\otimes\pi$ contain a single copy of
  $\pi _0$.
\end{itemize}

\end{frame}


\begin{frame}\frametitle{DHR program: derive structure of quantities}

  What is this unfamiliar mathematical structure? \pause

\bigskip
\begin{itemize}
\item A \emph{category} consists of objects (e.g.\ the
  representations), and arrows (e.g.\ transformations between
  representations). \pause
\item A \emph{tensor category} is a category with a binary operation
  $\otimes$ on objects and arrows, and a tensor unit
  $1_{\mathcal{C}}$, that satisfy certain laws (e.g.\ associativity). \pause
\item A tensor category has \emph{conjugates} just in case for each
  object $X$ there is an object $\ol X$ such that $X\otimes \ol X$ and
  $\ol X\otimes X$ contain a single copy of the tensor unit
  $1_{\mathcal{C}}$.
\end{itemize}
\end{frame}


\begin{frame}\frametitle{De-mystifying the magic symmetry principle}
\begin{thm}[Doplicher-Roberts 1990] The category of UIRs is equivalent to
  the dual of a symmetry group.  
\end{thm}

\end{frame}




\begin{frame}\frametitle{Metaphysical lessons}
\begin{itemize}
\item Stop equating ``quantity'' with ``something representable by
  real numbers'' --- that is just a \emph{special case}.  (Corollary:
  Rethink the contrast between ``qualitative'' and ``quantitative''.) \pause
\item Contra Mundy: don't need to resort to a second order theory of
  quantity, because any two physical objects have a ``sum''. \pause
\item Contra Mundy and Denby: determinates do have structure, but look
  to fundamental physics to tell us what it is. \pause
% \item Contra Balashov: having value zero for a quantity is not the
%  same as lacking a quantity. \pause
\item Contra the trend of relativizing: ``zero'' is in some cases
  objective and absolute.
\end{itemize}

\end{frame}

\end{document}
